\subsection{稀有事件的理想独立抽样基准}
\label{sec:rare-iid}
这里以正式任务表中的保留状态条目数说明预算量级，不改变原独立重复单位。L2在1024/4096预算分别有216/198份有效四链拟合，W1各有216份。设事件概率为$p$，假想$N$个状态均为独立后验样本，则命中数的期望为$Np$，零命中概率为$(1-p)^N$。表\ref{tab:rare-iid}使用两套L2重要性点估计及W1区间中点；L2参考有近似MCSE，W1中点只用于量级展示。
\begin{table}[htbp]\centering\small
\begin{tabular}{lrrrr}\toprule
参考 & 保留状态条目数 & $p$ & $Np$ & $(1-p)^N$\\\midrule
L2 Q1 & 4128768 & $1.56548\times10^{-9}$ & 0.0064635 & 99.356\%\\
L2 Q2 & 4128768 & $1.76048\times10^{-9}$ & 0.00726861 & 99.276\%\\
W1 & 4423680 & $4.76721\times10^{-11}$ & 0.000210886 & 99.979\%\\
\bottomrule\end{tabular}
\caption{全零事件的理想IID参照。实际MCMC相关，同核执行、设备和预算间还共享部分轨迹；这些条目不是独立信息数，表中概率不作为实际运行的零命中概率。}
\label{tab:rare-iid}\end{table}
\FloatBarrier
独立假设下，经验频率的相对标准误为$\sqrt{\frac{1-p}{Np}}$；当前即使使用表中的总条目数，仍远大于1。更一般地，只要每个状态的边缘分布正确，线性期望和并集界给出$\E[K]=Np$和$\Pr(K\geq1)\leq Np$，无需独立性；未平稳的实际初值不自动满足这一前提。全零结果本身因此不足以诊断混合不良，它揭示的是普通经验频率在当前预算下无法提供有用的稀有事件相对精度。常量事件的$\hat R$/ESS仍未定，与其他函数的探索诊断分开解释。
